% concatenated from Mehler_formula.tex
\documentclass[12pt]{article}
\usepackage{setspace} 
%\doublespacing
\usepackage[dvips]{graphicx} 
\usepackage{amsmath,amsthm,amsfonts,amssymb}
\RequirePackage[colorlinks,citecolor=blue,urlcolor=blue,linkcolor=blue]{hyperref}
\hypersetup{
colorlinks = true,
citecolor=blue,
urlcolor=blue,
linkcolor=blue,
pdfauthor = {Alexey Kuznetsov, Minjian Yuan},
pdfkeywords = {Wronskian, Hermite polynomials, Mehler formula, heat kernel, Darboux transformation},
pdftitle = {Mehler formula for Wronskians of Hermite polynomials},
pdfpagemode = UseNone
}
%****************************************************************************************************************
    \oddsidemargin -1.0cm
    \evensidemargin -1.0cm
    \topmargin -1.5cm
    \textwidth 18.7cm
    \textheight 23.5cm
%****************************************************************************************************************
    \def\qed{\hfill$\sqcap\kern-8.0pt\hbox{$\sqcup$}$\\}
    \def\q{{\mathbb Q}}
    \def\ind{{\mathbb I}}
    \def\re{\textnormal {Re}}
    \def\im{\textnormal {Im}}
    \def\p{{\mathbb P}}
    \def\e{{\mathbb E}}
    \def\E{{\mathcal E}}
    \def\r{{\mathbb R}}
    \def\wr{\textnormal{Wr}}
    \def\c{{\mathbb C}}
    \def\z{{\mathbb Z}} 
    \def\n{{\mathbb N}}  
    \def\d{{\textnormal d}}
    \def\i{{\textnormal i}} 
    \def\k{\textbf{\textsl{k}}}
    \def\j{\textbf{\textsl{j}}}    
%****************************************************************************************************************
	\newtheorem{theorem}{Theorem}
	\newtheorem{lemma}{Lemma}
	\newtheorem{proposition}{Proposition}
	\newtheorem{corollary}{Corollary}
	\newtheorem{definition}{Definition}
	\newtheorem{example}{Example}
	\newtheorem{exercise}{Exercise}
	\newtheorem{assumption}{Assumption}
	\newtheorem{remark}{Remark}

%****************************************************************************************************************
%****************************************************************************************************************
%****************************************************************************************************************

\title{Mehler formula for Wronskians of Hermite polynomials}
\author{ 
{Alexey Kuznetsov, Minjian Yuan \footnote{Dept. of Mathematics and Statistics,  York University,
4700 Keele Street, Toronto, ON, M3J 1P3, Canada.  
\\  Email: akuznets@yorku.ca, yuanm@yorku.ca}}}
 


\date{\today}


%**************************************************************************************************
%**************************************************************************************************
%**************************************************************************************************
\begin{document}
%**************************************************************************************************
%**************************************************************************************************
%**************************************************************************************************
\maketitle


\begin{abstract}
We prove that the bilinear generating function for Wronskians of Hermite polynomials can be expressed as the classical Mehler kernel multiplied by a polynomial, thereby extending the result of Pupasov-Maksimov \cite{Pupasov_2015} for exceptional Hermite polynomials. We establish several properties of the polynomials appearing in this extended version of the Mehler formula and present four conjectures about them.
\end{abstract}
%
{\vskip 0.15cm}
 \noindent {\it Keywords}: Wronskian, Hermite polynomials, Mehler formula, heat kernel, Darboux transformation
 
{\vskip 0.25cm}
 \noindent {\it 2020 Mathematics Subject Classification}: Primary 26C05, Secondary  33C45



%**************************************************************************************************
%**************************************************************************************************
%**************************************************************************************************



\section{Introduction and main results}

 Let $H_n(x)$ denote the Hermite polynomials. 
It is known (see \cite[formula (18.18.28)]{NIST}) that for all $|z|<1$ and $x,y\in \c$ 
\[
E(z,x,y):=\sum\limits_{n\geqslant 0} z^n \frac{H_n(x) H_n(y)}{2^n n!}=
\frac{1}{\sqrt{1-z^2}} \exp\bigg( \frac{2xzy-z^2(x^2+y^2)}{(1-z^2)} \bigg). 
\]
The expression on the right-hand side is called the Mehler kernel, and it plays a fundamental role in many applications. For example, it is used to construct the kernel of the fractional Fourier transform \cite{Namias1980}, the propagator for the harmonic oscillator in quantum mechanics \cite{Pupasov_2015}, and the transition probability density of the Ornstein-Uhlenbeck process \cite{Bakry2014}.  


Before stating our results, we introduce some notation. We denote by $\n$ the set of positive integers and by $\z_{\geqslant 0}$ the set of nonnegative integers. The Wronskian of smooth functions $\{f_j(x)\}_{1\leqslant j \leqslant m}$ is defined as 
\[
\wr[f_1,\dots,f_m]:=\det\,[\partial_x^{i-1} f_j(x)]_{1\leqslant i,j \leqslant m}.
\]  
For an $m$-tuple $\j=(j_1,j_2,\dots,j_m) \in \z_{\geqslant 0}^m$ we denote
\[
H_{\j}(x):=\wr[H_{j_1}(x), \dots,H_{j_m}(x)].
\]
We set $H_{\j}(x) \equiv 1$ if ${\mathbf j}=\varnothing$ (the empty tuple). 
In what follows, $\k=(k_1,k_2,\dots,k_l)\in \n^l$ will always denote a strictly increasing sequence of positive integers, and $l$ will denote  the length of the sequence $\k$. For $n\in \z_{\geqslant 0}$ we denote $H_{\k,n}(x):=H_{\j}(x)$  where $\j=(k_1,k_2,\dots,k_l,n)$. 

Wronskians of Hermite polynomials are important objects that arise in several contexts.  They play a fundamental role in the classification of rational monodromy-free Schr\"odinger operators with rational potentials having quadratic growth at infinity \cite{Oblomkov1999}. They appear as rational solutions of the fourth Painlev\'e equation \cite{Gomez2021}. Zeros of Wronskians of Hermite polynomials were studied in \cite{Felder2012}. It is known \cite{Adler1994} that $H_{\k}(x)$ has no real zeros if and only if $\k$ is a {\it Krein-Adler sequence}, that is, a strictly increasing sequence   $\k=(k_1,k_2,\dots, k_l) \in \n^l$ such that $l$ is even and  $k_{2i}=k_{2i-1}+1$ for all $1\leqslant i \leqslant l/2$.


We now define our main object of interest: 
\begin{equation}\label{def:Ek}
E_{\k}(z,x,y):=\sum\limits_{n \in \z_{\geqslant 0} \setminus \k}  z^n \frac{H_{\k,n}(x) H_{\k,n}(y)}{2^{n+l} n! \prod\limits_{j=1}^l (n-k_j)}.
\end{equation} 
As we will show later in Lemma \ref{lemma1}, the series in \eqref{def:Ek} converges for $|z|<1$, $x,y\in \c$ and defines an analytic function of three variables in that domain. Our main result is the following
\begin{theorem}\label{thm_main}
For any strictly increasing sequence $\k \in \n^l$ there exist polynomials $Q^{\k}_n(x,y)$, $0\leqslant n \leqslant k_l+1$ such that 
\begin{equation}\label{eqn_main}
E_{\k}(z,x,y)=E(z,x,y) \sum\limits_{n=0}^{k_l+1} z^n Q^{\k}_n(x,y),
\end{equation}
for all $|z|<1$, $x,y\in \c$.
\end{theorem}


Theorem \ref{thm_main} was established by A. M. Pupasov-Maksimov \cite{Pupasov_2015} in the important case where $\k$ is a  Krein-Adler sequence. To explain the significance of this result and the role that Krein-Adler sequences play, let us consider the functions  
\begin{equation}\label{def_fkn}
f_{\k,n}(x):=e^{-x^2/2} 
\frac{H_{\k,n}(x)}{H_{\k}(x)}. 
\end{equation}
It is known that for every strictly increasing sequence $\k \in \n^l$ these functions 
are solutions of the second-order linear differential equation 
\begin{equation}\label{f_kn_ODE}
{\mathcal L} f(x):=-f''(x)+U_{\k}(x) f(x)=(2n+1-2l) f(x),
\end{equation}
where $U_{\k}(x)$ is a rational function defined by
\begin{equation*}
U_{\k}(x):= x^2- 2\partial_x^2 \ln(H_{\k}(x))=
 x^2 - 2\frac{H_{\k}''(x)}{H_{\k}(x)}+
 2\frac{H_{\k}'(x)^2}{H_{\k}(x)^2}.
\end{equation*}
This  result follows from \cite[Proposition 5.2]{Gomez-Ullate_2014} and it is stated in this exact form in the proof of \cite[Proposition 5.5]{Gomez-Ullate_2014}. The differential operator in \eqref{f_kn_ODE} and the solutions of equation \eqref{f_kn_ODE} can be obtained by a sequence of Darboux transformations, see \cite[Theorem 3]{Oblomkov1999} and \cite{Gomez-Ullate_2014}.


As mentioned above, when $\k$ is a Krein-Adler sequence, the polynomial $H_{\k}(x)$ has no real zeros (see \cite[Theorem 4.2]{Gomez-Ullate_2014}) and the polynomials $H_{\k,n}(x)$ are called {\it exceptional Hermite polynomials}.  In this case, the functions $f_{\k,n}(x)$ belong to the Schwartz class (the class of smooth rapidly decreasing functions on $\r$) and satisfy the orthogonality condition
\begin{equation*}
\int_{\r} f_{\k,n}(x) f_{\k,m}(x) \d x= \delta_{n,m} \sqrt{\pi} 2^{n+l} n!  \prod_{j=1}^l (n-k_j), \;\;\; n,m \in {\mathbb Z}_{\geqslant 0} \setminus \k.
\end{equation*}
Moreover, the functions $\{f_{\k,n}(x)\}_{n \in {\mathbb Z}_{\geqslant 0} \setminus \k}$ form a complete orthogonal basis of $L_2(\r, \d x)$. These results can be found in \cite{Gomez-Ullate_2014,Gomez2016}. Thus, when $\k$ is a Krein-Adler sequence, the functions $\{f_{\k,n}(x)\}_{n \in {\mathbb Z}_{\geqslant 0} \setminus \k}$ form a complete eigenbasis for the operator ${\mathcal L}$ in $L_2(\r, \d x)$. This allows one to write down the spectral expansion of the corresponding heat kernel \cite{Avramidi2015,Vassilevich2003} (the integral kernel of the operator $\exp(-t {\mathcal L})$) in the form
\begin{equation}\label{L_heat_kernel}
p(t,x,y)=\sum\limits_{n \in {\mathbb Z}_{\geqslant 0} \setminus \k} 
e^{-(2n+1-2l) t} \; \frac{f_{\k,n}(x) f_{\k,n}(y)}{\lVert f_{\k,n} \rVert^2}=\frac{1}{\sqrt{\pi}} e^{(2l-1)t-(x^2+y^2)/2}
\;\frac{E_{\k}(e^{-2t},x,y)}{H_{\k}(x) H_{\k}(y)}. 
\end{equation}
Thus, the significance of Theorem \ref{thm_main} in the Krein-Adler case is that it gives an explicit expression for the heat kernel of the Schr\"odinger operator ${\mathcal L}$. 





\vspace{0.15cm}
\noindent
{\bf Remark 1:} Pupasov-Maksimov \cite{Pupasov_2015} states his results using the probabilists' version of Hermite polynomials, defined by ${\textnormal{He}}_n(x):=2^{-n/2} H_n(x/\sqrt{2})$. Denoting the polynomials appearing in \cite{Pupasov_2015} by $\widetilde Q_n^{\k}(x,y)$,  the relation between these two families of polynomials is  
\[
\widetilde Q_n^{\k}(x,y)=C \times  Q_n^{\k}(x/\sqrt{2}, y/\sqrt{2} ),
\]
where $C$ is a normalization constant that depends only on $\k$. 
\vspace{0.15cm}

In the following three propositions, we state several properties of the polynomials $Q_n^{\k}(x,y)$. 
 
 
\begin{proposition}\label{proposition1}
${}$
\begin{itemize}
\item[(i)]
For $n=0,1,\dots,k_l+1$ the polynomials $Q_n^{\k}(x,y)$ satisfy
\[
Q^{\k}_n(x,y)=Q_n^{\k}(y,x)=Q^{\k}_n(-x,-y)=(-1)^{\kappa+n} Q^{\k}_n(-x,y),
\] where
$\kappa:=k_1+k_2+\dots+k_l-l(l+1)/2$.
\item[(ii)]
We have 
\begin{equation}\label{eqn_Qk0}
Q^{\k}_0(x,y)=(-2)^{l} \bigg[ \prod\limits_{j=1}^l k_j \bigg] \times H_{\k-1}(x) H_{\k-1}(y), 
\end{equation}
where $\k-1:=(k_1-1,\dots,k_l-1)$.
\item[(iii)]
For $n=1,2,\dots,k_l+1$
the polynomials $Q^{\k}_n(x,y)$ can be computed recursively via 
\begin{equation}\label{Q_m_recursion}
Q^{\k}_n(x,y)=
{\mathbf 1}_{\{n \notin \k\}}
\frac{H_{\k,n}(x) H_{\k,n}(y)}{2^{n+l} n! \prod_{j=1}^l (n-k_j)}-
\sum\limits_{m=1}^{n}  \frac{H_{m}(x) H_{m}(y)}{2^{m} m!}
Q^{\k}_{n-m}(x,y).
\end{equation}
\end{itemize}
\end{proposition} 

The parity conditions and formula 
\eqref{Q_m_recursion} were established in \cite{Pupasov_2015} for Krein-Adler sequences $\k$. 



\begin{proposition}\label{proposition2}
Let $\k \in \n^l$ be a Krein-Adler sequence. Denote 
\begin{align*}
\Omega(x,y)&:=\frac{1}{y-x} \Big( \frac{H'_{\k}(y)}{H_{\k}(y)}
- \frac{H'_{\k}(x)}{H_{\k}(x)} \Big),
\\
\omega(x)&:=\lim_{y\to x} \Omega(x,y)=
 \frac{H_{\k}''(x)}{H_{\k}(x)}-\frac{H_{\k}'(x)^2}{H_{\k}(x)^2}.
\end{align*}
For every $x,y\in \c$, the following identities hold:
\begin{align}
\label{sum1}
\sum\limits_{n=0}^{k_l+1} Q^{\k}_n(x,y)&= H_{\k}(x) H_{\k}(y),\\
\label{sum2}
\sum\limits_{n=0}^{k_l+1} n Q_n^{\k}(x,y)&=- H_{\k}(x) 
H_{\k}(y)\big(  \Omega(x,y)-l \big),\\
\label{sum3}
\sum\limits_{n=0}^{k_l+1} n^2 Q_n^{\k}(x,y)&= H_{\k}(x) 
H_{\k}(y) \times \bigg((\Omega(x,y)-l)^2+\frac{\omega(x)+\omega(y)-2 \Omega(x,y)}{(y-x)^2}  \bigg).
\end{align}
\end{proposition}

Formula \eqref{sum1} for general $x$ and $y$ and formula \eqref{sum2} for the special case $x=y$ were first established in \cite{Pupasov_2015}.


\begin{proposition}\label{proposition3}
Let $k\in \n$ and $\k=(k)$, so that $l=1$. Then 
\begin{equation}\label{Q_n_l=1}
Q_n^{\k}(x,y)={\mathbf 1}_{\{n\geqslant 1\}} 2^{n-1} (k)_{n-1} H_{k+1-n}(x) H_{k+1-n}(y) - 2^{n+1} (k)_{n+1} 
H_{k-1-n}(x) H_{k-1-n}(y),
\end{equation}
where $(r)_n:=r(r-1)\dots(r-n+1)$ denotes the falling factorial. 
\end{proposition}

We also state several conjectures, which we have verified numerically for many cases of strictly increasing sequences $\k \in \n^l$.

\vspace{0.2cm}
\noindent
{\bf Conjecture 1:}  Identities \eqref{sum1}-\eqref{sum3} hold for any strictly increasing sequence $\k \in \n^l$, not only for Krein-Adler sequences. 


\vspace{0.2cm}
\noindent
{\bf Conjecture 2:} 
For $l\geqslant 2$ and any strictly increasing sequence $\k \in \n^l$, the following identity holds:
\begin{equation*}
Q^{\k}_{k_l+1} (x,y)=2^{k_l+l-1}
(k_l)!  \bigg[ \prod\limits_{j=1}^{l-1} (k_l-k_j) \bigg] \times H_{\tilde \k}(x) H_{\tilde \k}(y),
\end{equation*}
where $\tilde \k=(k_1,k_2,\dots,k_{l-1})$. 

\vspace{0.2cm}
\noindent
{\bf Conjecture 3:} When $l=2$ and $\k$ is a Krein-Adler sequence $\k=[k,k+1]$ (with $k \in \n$), then 
\begin{align*}
Q^{\k}_{k+1}(x,y)&=
2^{k+1} k! k \Big( H_{k+1}(x) H_{k+1}(y) 
+ 2 (k+1) \big[ H_{k+1}(x) H_{k-1}(y)\\
&\qquad \qquad +H_{k+1}(y) H_{k-1}(x) \big] + 4 k (k+1) H_{k-1}(x) H_{k-1}(y) \Big).
\end{align*}


\vspace{0.2cm}
\noindent
{\bf Conjecture 4:} 
The polynomials $Q^{\k}_n(x,y)$ have integer coefficients. A stronger conjecture is that the coefficients of the polynomials $Q^{\k}_n(x/2,y/2)$ are integers that are divisible by 
\[
 \prod\limits_{1\leqslant i < j \leqslant l} (k_i-k_j)^{2}.
 \]
 
Foata \cite{Foata1978} gave a combinatorial  proof of the classical Mehler formula. If true, Conjecture 4  may indicate that Theorem \ref{thm_main} could be proved by combinatorial methods and may suggest a potential combinatorial interpretation of the polynomials $Q_n^{\k}(x,y)$.

We present the proofs of all our results in the next section. 
Our proof of Theorem \ref{thm_main} follows the same general strategy as the proof of Pupasov-Maksimov \cite{Pupasov_2015} in the case of Krein-Adler sequences. However, the Krein-Adler condition, which implies the orthogonality and completeness of $f_{\k,n}$ in $L_2(\r, \d x)$, played an essential role in the proof in \cite{Pupasov_2015}. For the general case, we therefore had to modify the argument substantially
and introduce several new ideas.  

\section{Proofs}\label{section_proofs}

We denote by $\psi_n(x):=e^{-x^2/2} H_n(x)$ the non-normalized Hermite functions. If $n<0$  we set $H_n(x)\equiv 0$ and $\psi_n(x) \equiv 0$. For $\j=(j_1,j_2,\dots,j_m) \in \z_{\geqslant 0}^m$ we denote $\psi_{\j}(x):=\wr[\psi_{j_1}(x), \dots,\psi_{j_m}(x)]$ and
\[
|\,\j\,|:=j_1+\dots+j_m.
\]
For a strictly increasing sequence $\k \in \n^l$ and $n \in {\mathbb Z}_{\geqslant 0}$ we denote $\psi_{\k,n}(x):=\psi_{\j}(x)$, where $\j=(k_1,k_2,\dots,k_l,n)$. By the following well-known property of Wronskian determinants:
\begin{equation}\label{multiplication_identity}
\wr[f(x) g_1(x),f(x) g_2(x) \dots,f(x) g_{m}(x)]=f(x)^m \wr[ g_1(x), g_2(x) \dots, g_{m}(x)],
\end{equation}
we have 
\begin{equation}\label{psi_k_H_k}
\psi_{\k}(x)=e^{-l x^2/2} H_{\k}(x), \;\;\; \psi_{\k,n}(x)=e^{-(l+1) x^2/2} H_{\k,n}(x).
\end{equation}
In what follows, we work with the bilinear generating functions of $\psi_n$ and $\psi_{\k,n}$, defined by 
\begin{equation}\label{def_mathcal_E}
{\mathcal E}(z,x,y):=\sum\limits_{n\geqslant 0} z^n \frac{\psi_n(x) \psi_n(y)}{2^n n!}=
\frac{1}{\sqrt{1-z^2}} \exp\bigg( \frac{4xyz-(1+z^2)(x^2+y^2)}{2(1-z^2)} \bigg),
\end{equation}
and
\begin{equation}\label{def_mathcal_Ek}
{\mathcal E}_{\k}(z,x,y):=\sum\limits_{n \in \z_{\geqslant 0} \setminus \k}  z^n \frac{\psi_{\k,n}(x) \psi_{\k,n}(y)}{2^{n+l} n! \prod\limits_{j=1}^l (n-k_j)}.
\end{equation} 
From \eqref{psi_k_H_k}, we have 
\begin{equation}\label{mathcal_E_E}
{\mathcal E}(z,x,y)=e^{-(x^2+y^2)/2} E(z,x,y), \;\;\; 
{\mathcal E}_{\k}(z,x,y)=e^{-(l+1)(x^2+y^2)/2} E_{\k}(z,x,y).
\end{equation}




 
Before proving Theorem \ref{thm_main}, we need to establish several auxiliary results.  


\begin{lemma}\label{lemma1}
${}$
\begin{itemize}
\item[(i)]
There exists $C=C(\k)>0$ such that for all $n\geqslant 1$ and $x\in \c$
\begin{equation}\label{psi_upper_bound}
|\psi_{\k,n}(x)|\leqslant C n^l \big | e^{-l x^2/2}\big | \times (1+|x|^{|\k|}) \times \sum\limits_{j=0}^l 
|\psi_{n-j}(x)|.
\end{equation}
\item[(ii)]
 The function $\E_{\k}(z,x,y)$ 
is an analytic function of three variables in the domain $|z|<1$, $x,y\in \c$. 
\item[(iii)] The function $\E_{\k}(z,x,y)$  satisfies the following symmetry conditions:
\begin{equation}\label{E_k_symmetries}
\E_{\k}(z,x,y)=\E_{\k}(z,y,x)=\E_{\k}(z,-x,-y)=(-1)^{\kappa} \E_{\k}(-z,-x,y),  \;\;\; |z|<1, \; x,y\in \c, 
\end{equation}
where $\kappa=|\,\k\,|-l(l+1)/2$.
\end{itemize} 
\end{lemma}
\begin{proof}
To prove (i), we start with the formula 
\[
\psi_{\k,n}(x)=e^{-(l+1)x^2/2} \wr[H_{k_1}(x),\dots,H_{k_l}(x),H_n(x)],
\]
and expand the Wronskian determinant along the last column, which contains the derivatives of $H_n(x)$:
\begin{equation}\label{psi_k_n_sum_derivatives}
\psi_{\k,n}(x)=e^{-(l+1)x^2/2} \sum\limits_{j=0}^l p_{j}(x) \partial_x^j H_n(x),
\end{equation}
where
\[
p_{j}(x)=(-1)^{j+l} {\textnormal{det}}
\Big[ \partial_x^m H_{k_i}(x) \Big]_{\substack{1\leqslant i \leqslant l \\ 0\leqslant m \leqslant l \\ m\neq j}}.
\]
It is clear from the above determinant expression that $p_{j}$ is a polynomial of degree not exceeding $|\, \k \, |$. This, combined with the formula 
\[
\partial_x^j H_m(x)=2^j m(m-1)\dots(m-j+1) H_{m-j}(x),
\]
implies the upper bound \eqref{psi_upper_bound}. 


The proof of item (ii) requires the following asymptotic result (see \cite{Szego}[Theorem 8.22.7]):  as  $n\to +\infty$
\[
\psi_n(x) = \frac{2^n}{\sqrt{\pi}} \Gamma((n+1)/2) \Big[ \cos(x \sqrt{2n+1} - n\pi/2)+O\big(n^{-1/2} \exp\big( |\im(x)|\sqrt{2n+1}\big)\big)\Big], 
\]
uniformly in $x$ on compact subsets of $\c$. We also need the fact that 
\[
\frac{\Gamma((n+1)/2)^2}{n!}=2^{-n} \sqrt{\frac{2\pi}{n}}  (1+o(1)), \;\;\; n\to +\infty,
\]
which follows from the Legendre duplication formula for the gamma function. 
Combining the above two facts with \eqref{psi_upper_bound}, we conclude that for each compact subset $A \subset \c$ there exists a constant $C_1=C_1(A,\k)>0$ such that the absolute value of each term in the infinite series \eqref{def_mathcal_Ek} is bounded above by 
\[
C_1 |z|^n n^{l-1/2} \exp\big(\sqrt{2n+1} (|\im(x)|+|\im(y)|)\big), 
\]
for $x,y \in A$. This implies that the series \eqref{def_mathcal_Ek} converges uniformly on compact subsets of $\{(z,x,y) \in \c^3 \; : \; |z|<1\}$
and defines an analytic function in this domain. 


According to Lemma 2.1 and Lemma 3.6 in \cite{Bonneux2018}, the polynomial $H_{\k,n}(x)$ has degree $\kappa+n$ and satisfies the parity condition $H_{\k,n}(-x)=(-1)^{\kappa+n} H_{\k,n}(x)$. 
The symmetry conditions in item (iii) follow immediately from this result and formulas \eqref{psi_k_H_k}, \eqref{def_mathcal_E} and \eqref{def_mathcal_Ek}.
\end{proof}





\begin{lemma}\label{lemma2}
Assume that $U(x)$, $g_1(x)$ and $g_2(x)$ are smooth functions on $(-\infty, c)$ and that
\begin{equation}\label{f_i_equation}
 -g_i''(x)+U(x) g_i(x)=\lambda_i g_i(x), \;\;\;  x \in (-\infty, c), \;\;\; i\in \{1,2\}. 
\end{equation}
 Assume further that $\lambda_1\neq \lambda_2$ and $\wr[g_1(x),g_2(x)] \to 0$ as $x\to -\infty$. Then, for any $x\in (-\infty, c)$, we have
\[
\int_{-\infty}^x g_1(y)g_2(y) \d y = \frac{1}{ \lambda_1-\lambda_2} \wr[g_1(x),g_2(x)].
\]
\end{lemma}
\begin{proof} The proof is based on the identity
\[
\frac{\d}{\d x}\wr[g_1(x),g_2(x)] =  (\lambda_1-\lambda_2) g_1(x)g_2(x),
\]
which follows from \eqref{f_i_equation}.
\end{proof}




 

Let ${\mathcal L}_{\k,x}$ be the first-order  differential operator acting on the $x$-variable, defined by
\[
{\mathcal L}_{\k,x}h(x):=\wr[ \psi_{\k}(x), h(x)]=\psi_{\k}(x) h'(x)
-\psi_{\k}'(x) h(x),
\]
where $h$ is a smooth function on $\r$. 
For a strictly increasing sequence $k=(k_1,\dots,k_l)\in \n^l$ we denote by $\k \setminus k_j$ the sequence of length $l-1$ obtained from $\k$ by removing the term $k_j$.  We also denote 
\[
\tilde \k:=\k \setminus k_l=(k_1,\dots,k_{l-1}).
\] The Wronskian identity
\begin{equation}\label{Wronskian_identity}
\wr[g_1,\dots,g_{j}]\times \wr[g_1,\dots,g_{j},h,f] 
=\wr\Big[\wr[g_1,\dots,g_{j},h],\wr[g_1,\dots,g_{j},f] \Big]
\end{equation}
 implies the following key result
\begin{equation}\label{eqn_psi_kn_L_k}
\psi_{\tilde \k}(x) \times \psi_{\k,n}(x)= {\mathcal L}_{\k,x} \psi_{\tilde \k,n}(x).
\end{equation}


\begin{lemma}\label{lemma3} Let $\k \in \n^l$ be a strictly increasing sequence. There exists $c=c(\k) \in \r$ such that for $|z|<1$, $x\in \r$ and $y \in (-\infty,c]$ we have
\begin{equation}\label{eqn_E_recursion}
\E_{\k}(z,x,y)=-  \frac{\psi_{\tilde \k}(y)}{\psi_{\tilde \k}(x)} 
{\mathcal L}_{\k,x} \int_{-\infty}^y 
\E_{\tilde \k}(z,x,u) \frac{\psi_{\k}(u)}{\psi_{\tilde \k}(u)^2} \d u.
\end{equation}
\end{lemma}
\begin{proof}
We take $c \in \r$ such that the polynomial $H_{\tilde \k}(x)$ has no zeros on $(-\infty,c]$. This implies that $\psi_{\tilde \k}(u)=\exp(-(l-1)u^2/2) H_{\tilde \k}(u)$ is also non-zero on $(-\infty,c]$. We recall that $f_{\k,n}(x)$ is defined in \eqref{def_fkn} and that these functions satisfy  equation \eqref{f_kn_ODE}.  

To prove formula \eqref{eqn_E_recursion}, we first apply Lemma \ref{lemma2}:
\begin{align}\label{lemma3_proof1}
&\int_{-\infty}^y \E_{\tilde \k}(z,x,u) \frac{\psi_{\k}(u)}{\psi_{\tilde \k}(u)^2} \d u=
\sum\limits_{n \in \z_{\geqslant 0} \setminus \tilde \k} z^n
\frac{\psi_{\tilde \k,n}(x)}{2^{n+l-1} n! \prod\limits_{j=1}^{l-1} (n-k_j)} \int_{-\infty}^y f_{\tilde \k,k_l}(u) f_{\tilde \k,n}(u) \d u\\ \nonumber
&=\sum\limits_{n \in \z_{\geqslant 0} \setminus \k} z^n
\frac{ \psi_{\tilde \k,n}(x)}{2^{n+l-1} n! \prod\limits_{j=1}^{l-1} (n-k_j)} 
\frac{\wr[f_{\tilde\k,k_l}(y),f_{\tilde \k,n}(y)]}{2k_l-2n}+ 
\frac{z^{k_l} \psi_{\tilde \k,k_l}(x)}{2^{k_l+l-1} k_l! \prod\limits_{j=1}^{l-1} (k_l-k_j)}  \int_{-\infty}^y f_{\tilde \k,k_l}(u)^2  \d u.
\end{align}
From \eqref{eqn_psi_kn_L_k} we find
\[
\wr[f_{\tilde\k,k_l}(y),f_{\tilde \k,n}(y)]=\frac{1}{\psi_{\tilde \k}(y)^2} \wr[\psi_{\k}(y),\psi_{\tilde \k,n}(y)]=
\frac{\psi_{\k,n}(y)}{\psi_{\tilde \k}(y)},
\]
and therefore
\[
\int_{-\infty}^y \E_{\tilde \k}(z,x,u) \frac{\psi_{\k}(u)}{\psi_{\tilde \k}(u)^2} \d u=
-\frac{1}{\psi_{\tilde \k}(y)}\sum\limits_{n \in \z_{\geqslant 0} \setminus \k} z^n
\frac{\psi_{\tilde \k,n}(x) \psi_{\k,n}(y)}{2^{n+l} n! \prod\limits_{j=1}^{l} (n-k_j)}+z^{k_l} \psi_{\k}(x) G(y),
\]
for some function $G(y)$. 
Finally, applying identity \eqref{eqn_psi_kn_L_k} once again and noting that ${\mathcal L}_{\k,x} \psi_{\k}(x)\equiv 0$, we obtain
\begin{align}\label{lemma3_proof2}
-  \frac{\psi_{\tilde \k}(y)}{\psi_{\tilde \k}(x)} 
{\mathcal L}_{\k,x} \int_{-\infty}^y 
\E_{\tilde \k}(z,x,u) \frac{\psi_{\k}(u)}{\psi_{\tilde \k}(u)^2} \d u&=
\frac{1}{\psi_{\tilde \k}(x)}\sum\limits_{n \in \z_{\geqslant 0} \setminus \k} z^n
\frac{{\mathcal L}_{\k,x} \psi_{\tilde \k,n}(x) \psi_{\k,n}(y)}{2^{n+l} n! \prod\limits_{j=1}^{l} (n-k_j)}
\\ \nonumber &=
\sum\limits_{n \in \z_{\geqslant 0} \setminus \k} z^n
\frac{\psi_{\k,n}(x) \psi_{\k,n}(y)}{2^{n+l} n! \prod\limits_{j=1}^{l} (n-k_j)}=\E_{\k}(z,x,y).
\end{align}

In the above derivation of formula \eqref{eqn_E_recursion},  we interchanged summation and differentiation/integration. Let us first justify the interchange of summation and integration in the first step of \eqref{lemma3_proof1}. Consider the function $f_{\tilde \k,k_l}(u) f_{\tilde \k,n}(u)$, which we rewrite in the form 
\[
f_{\tilde \k,k_l}(u) f_{\tilde \k,n}(u)=\Big[ e^{-u^2/2} \frac{H_{\k}(u)}{H_{\tilde \k}(u)^2} \Big] \times \Big[ e^{(l-1)u^2/2} \psi_{\tilde \k,n}(u)\Big].
\]
The above factorization follows from \eqref{def_fkn} and \eqref{psi_k_H_k}.
The function $e^{-u^2/4} H_{\k}(u)/H_{\tilde \k}(u)^2$ is bounded by some constant $A=A(\k)$ on $(-\infty, c]$ (since $H_{\k}(u)$ and $H_{\tilde \k}(u)$ are polynomials and $H_{\tilde \k}(u)$ is non-zero on $(-\infty, c]$). Thus,
$g(u):=e^{-u^2/2} H_{\k}(u)/H_{\tilde \k}(u)^2$ satisfies $|g(u)|<A e^{-u^2/4}$ for $u \in (-\infty, c]$.  
From \eqref{psi_upper_bound} we conclude that there exists a constant $B=B(\tilde \k)$ such that 
\[
|e^{(l-1)u^2/2} \psi_{\tilde \k,n}(u)|< B n^{l-1} (1+|u|^{M}) \sum\limits_{j=0}^{l-1} 
|\psi_{n-j}(u)|, \;\;\; u \in \r,
\]
where $M=|\tilde \k|$. 
Applying the Cauchy-Schwarz inequality and using the fact that  $\lVert \psi_n \rVert^2=\sqrt{\pi}2^n n!$, we obtain 
\begin{align*}
 \int_{-\infty}^y  | f_{\tilde \k,k_l}(u) f_{\tilde \k,n}(u) | \d u 
&\leqslant 
 A B n^{l-1} \sum\limits_{j=0}^{l-1}  \int_{-\infty}^c (1+|u|^{M}) e^{-u^2/4} \;
|\psi_{n-j}(u)|\, \d u
\\ &\leqslant A B  n^{l-1} \Big[\int_{-\infty}^c (1+|u|^{M})^2 e^{-u^2/2}\d u\Big]^{1/2} 
\times  \sum\limits_{j=0}^{l-1} \lVert \psi_{n-j}\rVert
\leqslant C l n^{l-1} \sqrt{2^n n!},
\end{align*}
for some $C=C(\k)>0$. 
Using the above estimate and the same argument as in the proof of Lemma \ref{lemma1}, we conclude that the series 
\[
\sum\limits_{n \in \z_{\geqslant 0} \setminus  \k} r^n
\frac{ |\psi_{\tilde \k,n}(x)|}{2^{n+l-1} n! \prod\limits_{j=1}^{l-1} (n-k_j)}  \int_{-\infty}^y | f_{\tilde \k,k_l}(u) f_{\tilde \k,n}(u) | \d u 
 \]
 converges for all $r\in (0,1)$, $x\in \c$ and $y \in (-\infty, c]$. Thus, we can apply Fubini's Theorem and interchange summation and integration in \eqref{lemma3_proof1}.
 
To justify the interchange of summation and differentiation in the first step of \eqref{lemma3_proof2}, we use the same argument as in the proof of Lemma \ref{lemma1} and conclude that the series  
\[
\sum\limits_{n \in \z_{\geqslant 0} \setminus \k} z^n
\frac{\psi_{\tilde \k,n}(x) \psi_{\k,n}(y)}{2^{n+l} n! \prod\limits_{j=1}^{l} (n-k_j)}
\]
converges uniformly on compact subsets of $\{(z,x,y) \in \c^3 \; : \; |z|<1\}$. Hence,  it defines an analytic function on this set, and the infinite series can be differentiated term-by-term. 
\end{proof}


For $n\in \n$, $|z|<1$, $x\in \c$ and $y\in \r$ we denote 
\begin{equation}\label{def_In}
I_n(z,x,y):=\int_{-\infty}^y \E(z,x,u) \psi_n(u) \d u.
\end{equation}


\begin{lemma}\label{lemma4}
The function $I_n(z,x,y)$ extends to an analytic function in the domain \mbox{$\{(z,x,y) \in \c^3 \; : \; |z|<1\}$} and it satisfies 
\begin{equation}\label{eqn_I_n_main}
I_n(z,x,y)=
 z^n \psi_n(x)  \frac{\sqrt{\pi}}{2} \big( 1+ 
{\textnormal{erf}}(w) \big)-e^{-x^2/2-w^2} \sum\limits_{i=1}^n \binom{n}{i} \zeta^{i} z^{n-i}
H_{i-1}(w) H_{n-i}(x),
\end{equation}
where $\zeta:=\sqrt{1-z^2}$ and $w:=(y-xz)/\zeta$. 
\end{lemma}
\begin{proof}
We first observe that $\E(z,x,y)$ defined in \eqref{def_mathcal_E} can be written in the form 
\begin{equation*}
\E(z,x,y)=\zeta^{-1} e^{y^2/2-x^2/2-(y-xz)^2/\zeta^2}.
\end{equation*}
Assume that $z\in (-1,1)$ and $x,y\in \r$. 
Using the above expression for $\E(z,x,y)$, we rewrite \eqref{def_In} in the equivalent form 
\[
I_n(z,x,y)=\zeta^{-1} e^{-x^2/2} \int_{-\infty}^y e^{-(u-xz)^2/\zeta^2} H_n(u) \d u. 
\]
We change the variable of integration to $v=(u-xz)/\zeta$ and obtain
\begin{equation*}
I_n(z,x,y)=e^{-x^2/2} \int_{-\infty}^w e^{-v^2} H_n(\zeta v + x z) \d v. 
\end{equation*}
Noting that $\zeta^2+z^2=1$, we apply the addition formula for Hermite polynomials (see \cite[formula 18.18.11]{NIST})
\begin{equation*}
H_n(\zeta v+xz)=\sum\limits_{i=0}^n \binom{n}{i} \zeta^i z^{n-i} 
H_i(v) H_{n-i}(x).
\end{equation*}
The desired result \eqref{eqn_I_n_main} follows from the above two equations and the following formulas:
\begin{equation}\label{two_integrals}
\int_{-\infty}^w e^{-v^2} \d v= \frac{\sqrt{\pi}}{2} \big( 1+ 
{\textnormal{erf}}(w) \big), \;\;\; \int_{-\infty}^w e^{-v^2} H_i(v) \d v=-e^{-w^2} H_{i-1}(w),
\end{equation}
which can be found in \cite[\S 7.2]{NIST}  and \cite[eqn. (7.373)]{Jeffrey2007}. 

Now that we have established \eqref{eqn_I_n_main} for $z \in (-1,1)$, $x,y\in \r$, we note that the right-hand side in \eqref{eqn_I_n_main} is an analytic function in $\{(z,x,y) \in \c^3 \; : \; |z|<1\}$. Thus, $I_n(z,x,y)$ can be extended in this domain by analytic continuation. 
\end{proof}


We introduce the $l$-th order linear differential operator 
${\mathcal M}_{\k,x}$, acting on the $x$-variable, by
\begin{equation}\label{def_M_kx}
{\mathcal M}_{\k,x} h(x):=\wr[ \psi_{k_1},\dots,\psi_{k_l},h](x), 
\end{equation}
where $h$ is an entire function. The following result was the starting point for the derivation of the Mehler formula for exceptional Hermite polynomials in  \cite{Pupasov_2015}; see also \cite{Pupasov_2007}, where this result was first obtained for Krein-Adler sequences $\k$.   



\begin{lemma}\label{lemma5}
For $|z|<1$, $x,y\in \c$  we have
\begin{equation}\label{eqn_PM}
\E_{\k}(z,x,y)=  
{\mathcal M}_{\k,x} \sum\limits_{j=1}^l (-1)^{j+l-1} \psi_{\k \setminus k_j}(y)  I_{k_j}(z,x,y).
\end{equation}
\end{lemma}
\begin{proof}
The proof proceeds by induction on $l$. The case $l=1$ and $\k=(k)$ is covered by Lemma \ref{lemma3}, followed by analytic continuation in $x$ and $y$. 
 Assume now  that $l\geqslant 2$ and that for $\tilde \k=(k_1,k_2,\dots,k_{l-1})$ we have 
\[
\E_{\tilde \k}(z,x,y)= 
{\mathcal M}_{\tilde \k,x} \sum\limits_{j=1}^{l-1} (-1)^{j+l-2} \psi_{\tilde \k \setminus k_j}(y)  I_{k_j}(z,x,y),
\]
with the above identity valid for all $z \in (-1,1)$ and $x,y\in \r$.
We substitute the above expression for $\E_{\tilde \k}(z,x,y)$ into the right-hand side of \eqref{eqn_E_recursion} and conclude that there exists $c \in \r$ such that for $z\in (-1,1)$, $x\in \r$ and $y \in (-\infty,c]$ we have  
\begin{equation}\label{lemma5_proof1}
\E_{\k}(z,x,y)=
  \frac{\psi_{\tilde \k}(y)}{\psi_{\tilde \k}(x)} 
{\mathcal L}_{\k,x}{\mathcal M}_{\tilde \k,x} \sum\limits_{j=1}^{l-1} (-1)^{j+l-1} \int_{-\infty}^y 
 \psi_{\tilde \k \setminus k_j}(u)  I_{k_j}(z,x,u) \frac{\psi_{\k}(u)}{\psi_{\tilde \k}(u)^2} \d u.
\end{equation}
Using the Wronskian identity \eqref{Wronskian_identity}, we check that 
\[
\frac{\d}{\d u} \frac{\psi_{\k \setminus k_j}(u)}{\psi_{\tilde \k}(u)^2}=\psi_{\tilde \k \setminus k_j}(u) \frac{\psi_{\k}(u)}{\psi_{\tilde \k}(u)^2}.
\]
Note that $\psi_{\k \setminus k_j}(u)/\psi_{\tilde \k}(u)$ is a rational function of $u$ (this follows from \eqref{multiplication_identity}). Due to  our assumption $z\in (-1,1)$, we have $\zeta>0$, thus from \eqref{eqn_I_n_main} we see that the function $u\mapsto I_{k_j}(z,x,u)$ decays exponentially fast as $u\to -\infty$ (this fact also follows from the integral definition \eqref{def_In}). Thus, 
\[
\lim\limits_{u\to -\infty} I_{k_j}(z,x,u)\frac{\psi_{\k \setminus k_j}(u)}{\psi_{\tilde \k}(u)}=0.
\]
Using the above two results and integrating by parts, we obtain 
\begin{equation}\label{lemma5_proof2}
\int_{-\infty}^y 
 \psi_{\tilde \k \setminus k_j}(u)  I_{k_j}(z,x,u) \frac{\psi_{\k}(u)}{\psi_{\tilde \k}(u)^2} \d u=
 I_{k_j}(z,x,y) \frac{\psi_{\k \setminus k_j}(y)}{\psi_{\tilde \k}(y)}-\int_{-\infty}^y \E(z,x,u) \psi_{k_j}(u) 
 \frac{\psi_{\k \setminus k_j}(u)}{\psi_{\tilde \k}(u)} \d u.
\end{equation}


Next,  we use the identity
\[
0=\det \begin{bmatrix}
\psi_{k_1}(u) & \psi_{k_2}(u) & \cdots & \psi_{k_{l}}(u)\\
\psi_{k_1}(u) & \psi_{k_2}(u) & \cdots & \psi_{k_{l}}(u)\\
\psi_{k_1}'(u) & \psi_{k_2}'(u) & \cdots & \psi_{k_{l}}'(u)\\
\vdots & \vdots & \ddots & \vdots\\
\psi_{k_1}^{(l-2)}(u) & \psi_{k_2}^{(l-2)}(u) & \cdots & \psi_{k_{l}}^{(l-2)}(u)
\end{bmatrix}=\sum\limits_{j=1}^l (-1)^{j-1} \psi_{k_j}(u)\times  \psi_{\k \setminus k_j}(u), 
\]
where we have expanded the above determinant along the first row. 
Therefore, 
\[
\sum\limits_{j=1}^{l-1} (-1)^{j-1} \psi_{k_j}(u)  \psi_{\k \setminus k_j}(u) = (-1)^{l} \psi_{k_l}(u)  \psi_{\tilde \k}(u).
\]
Combining the above result with \eqref{lemma5_proof1} and \eqref{lemma5_proof2}
we arrive at 
\[
\E_{\k}(z,x,y)=
  \frac{1}{\psi_{\tilde \k}(x)} 
{\mathcal L}_{\k,x}{\mathcal M}_{\tilde \k,x} \sum\limits_{j=1}^{l} (-1)^{j+l-1} 
 \psi_{\k \setminus k_j}(y)  I_{k_j}(z,x,y). 
\]
From the Wronskian identity \eqref{Wronskian_identity} it follows that
\[
{\mathcal L}_{\k,x}{\mathcal M}_{\tilde \k,x} h(x) = \psi_{\tilde \k}(x) {\mathcal M}_{\k,x}h(x), 
\]
for any entire function $h$. Thus, we have proved that for $z\in (-1,1)$, $x\in \r$ and $y\in (-\infty, c]$ we have
\[
\E_{\k}(z,x,y)=
 {\mathcal M}_{\k,x} \sum\limits_{j=1}^{l} (-1)^{j+l-1} 
 \psi_{\k \setminus k_j}(y)  I_{k_j}(z,x,y).
\]
By analytic continuation, the above identity holds everywhere in the domain $\{(z,x,y)\in \c^3 \; : \; |z|<1\}$. 
This completes the induction step. 
\end{proof}


\begin{corollary}\label{Corollary1}
For $l\geqslant 2$,  $|z|<1$, $x\in \c$ and $y\in \r$  we have
\begin{equation}\label{eqn_PM2}
\E_{\k}(z,x,y)=  
{\mathcal M}_{\k,x}  \int_{-\infty}^y  \E(z,x,u) e^{-u^2/2} P_{\k}(u,y) \d u,
\end{equation}
where $P_{\k}(u,y)$ is a polynomial in $u$ that satisfies $\partial_u^j P_{\k}(u,y) |_{u=y}=0$ for all $j=0,1,\dots,l-2$. 
\end{corollary}
\begin{proof}
Formula \eqref{eqn_PM2} is equivalent to \eqref{eqn_PM} with $P_{\k}(u,y)$ defined by
\[
e^{-u^2/2} P_{\k}(u,y)=\sum\limits_{j=1}^l (-1)^{j+l-1}  \psi_{k_j}(u) \psi_{\k \setminus k_j}(y)=(-1)^l \det \begin{bmatrix}
\psi_{k_1}(u) & \psi_{k_2}(u) & \cdots & \psi_{k_{l}}(u)\\
\psi_{k_1}(y) & \psi_{k_2}(y) & \cdots & \psi_{k_{l}}(y)\\
\psi_{k_1}'(y) & \psi_{k_2}'(y) & \cdots & \psi_{k_{l}}'(y)\\
\vdots & \vdots & \ddots & \vdots\\
\psi_{k_1}^{(l-2)}(y) & \psi_{k_2}^{(l-2)}(y) & \cdots & \psi_{k_{l}}^{(l-2)}(y)
\end{bmatrix}.
\]
The above equation can be written in the form
\[
P_{\k}(u,y)=e^{-(l-1)y^2/2} (-1)^l \det\begin{bmatrix}
H_{k_1}(u) & H_{k_2}(u) & \cdots & H_{k_{l}}(u)\\
H_{k_1}(y) & H_{k_2}(y) & \cdots & H_{k_{l}}(y)\\
H_{k_1}'(y) & H_{k_2}'(y) & \cdots & H_{k_{l}}'(y)\\
\vdots & \vdots & \ddots & \vdots\\
H_{k_1}^{(l-2)}(y) & H_{k_2}^{(l-2)}(y) & \cdots & H_{k_{l}}^{(l-2)}(y)
\end{bmatrix},
\]
which immediately implies 
$\partial_u^j P_{\k}(u,y) |_{u=y}=0$ for all $j=0,1,\dots,l-2$. 
\end{proof}



\begin{lemma}\label{lemma6}
Let $m,n \in {\mathbb Z}_{\geqslant 0}$ and let $x, y$ be fixed real numbers satisfying $x>y$. Let $P(u)$ be a polynomial such that $P^{(i)}(y)=0$ for $i=0,1,\dots,n$. If $m\leqslant n+1$ then 
\begin{equation}\label{eqn_integral_limit}
\lim\limits_{z\to 1-} \frac{\partial_x^{m} \int_{-\infty}^y \E(z,x,u) e^{-u^2/2}  P(u) \d u}{\E(z,x,y)} 
=0. 
\end{equation}
If $m=n+2$, the above limit exists and is finite.  
\end{lemma}
\begin{proof}
First, we check that formula \eqref{def_mathcal_E} can be rewritten in the equivalent form  
\begin{equation}\label{E_factorization}
\E(z,x,y)  =e^{(x^2+y^2)/2  - 2xy/(1+z)} \times 
\zeta^{-1} 
e^{-(y-x)^2/\zeta^2 },
\end{equation}
where $\zeta=\zeta(z)=\sqrt{1-z^2}$. For the rest of this proof, we assume that $1/\sqrt{2}<z<1$, so that $0<\zeta<1/\sqrt{2}$.  


 
We claim that, for any $m\in {\mathbb Z}_{\geqslant 0}$ and any fixed $x,y \in \r$  such that $x>y$, we have
\begin{equation}\label{eqn_q_integral}
\zeta^{-2} e^{(y-x)^2/\zeta^2} \int_{-\infty}^y  
e^{-2xu/(1+z)-(u-x)^2/\zeta^2}   u^m \d u \to 
\frac{e^{-x y}y^m}{2(x-y)},  
\end{equation}
as $z\to 1-$ (note that $\zeta\to 0+$ as $z\to 1-$). 
Indeed, changing the variable of integration $u=y-\zeta^2 v$, we obtain
\[
\zeta^{-2}  e^{(y-x)^2/\zeta^2}  \int_{-\infty}^y  
e^{-2xu/(1+z)-(u-x)^2/\zeta^2}   u^m \d u
= e^{-2x y/(1+z)} \int_0^{\infty} e^{-2(x-y)v -\zeta^2 (v^2-2xv/(1+z))}  (y-\zeta^2 v)^m \d v,
\]
and the limit in \eqref{eqn_q_integral} follows by the dominated convergence theorem. 




Next, let $q(z,u)$ be a polynomial in $u$ of the form
\begin{equation}\label{form_of_q}
q(z,u)=\sum\limits_{i=0}^M a_i(z) u^i,
\end{equation}
where $M$ does not depend on $z$ and the coefficients $a_i(z)$ are continuous in some neighborhood of $z=1$. For $z\in (0,1)$ and $j \in {\mathbb Z}_{\geqslant 0}$ we consider 
\begin{equation}\label{def_Q_j}
{\mathcal Q}_{j}(z,x,y):=\zeta^{-1} \int_{-\infty}^y \Big[\partial_u^j  
e^{-(u-x)^2/\zeta^2} \Big] \times  e^{-2xu/(1+z)} q(z,u) (u-y)^{n+1} \d u. 
\end{equation}
From \eqref{E_factorization} and \eqref{eqn_q_integral} we conclude that 
\begin{equation*}
\frac{{\mathcal Q}_{0}(z,x,y)}{\E(z,x,y)} \to 0
\end{equation*}
as $z \to 1-$. When $0< j \leqslant n+1$, we 
integrate by parts $j$ times 
and obtain 
\begin{align}\label{Q_j_integration_by_parts}
{\mathcal Q}_j(z,x,y)&= 
\zeta^{-1} \sum\limits_{i=0}^{j-1} (-1)^i \Big[\partial_u^{j-1-i}  
e^{-(u-x)^2/\zeta^2} \Big]_{u=y} \times 
\Big[\partial_u^i \Big(  
e^{-2xu/(1+z)} q(z,u) (u-y)^{n+1}\Big)\Big]_{u=y} 
\\\nonumber &+(-1)^j \zeta^{-1} \int_{-\infty}^y 
e^{-(u-x)^2/\zeta^2} 
\partial_u^j  \Big(
e^{-2xu/(1+z)} q(z,u) (u-y)^{n+1}\Big) \d u.
\end{align}
In the finite sum in the above formula, we have $i\leqslant j-1 \leqslant n$, and therefore
\[
\Big[\partial_u^i  
e^{-2xu/(1+z)} q(z,u) (u-y)^{n+1}\Big]_{u=y}=0. 
\]
Thus, the finite sum in \eqref{Q_j_integration_by_parts} is equal to zero, and
it follows from \eqref{E_factorization} and \eqref{eqn_q_integral} that, for $0<j\leqslant n+1$,
\begin{equation*}
\frac{{\mathcal Q}_{j}(z,x,y)}{\E(z,x,y)} \to 0, 
\end{equation*}
as $z \to 1-$. 

When $j=n+2$, the preceding argument shows that  the terms with $0\leqslant i \leqslant j-2$ in the finite sum in \eqref{Q_j_integration_by_parts} are all equal to zero, and when $i=j-1=n+1$ we have
\[
\Big[\partial_u^{n+1} \Big(  
e^{-2xu/(1+z)} q(z,u) (u-y)^{n+1}\Big)\Big]_{u=y}= 
e^{-2xy/(1+z)} q(z,y) (n+1)!.
\]
Thus, it follows from 
\eqref{E_factorization}, \eqref{eqn_q_integral} and 
\eqref{Q_j_integration_by_parts} that, as $z\to 1-$,
\begin{equation*}
\frac{{\mathcal Q}_{n+2}(z,x,y)}{\E(z,x,y)} \to 
 (-1)^{n+1} 
 e^{-(x^2+y^2)/2} q(1,y) (n+1)!. 
\end{equation*}


Now we are ready to complete the proof of Lemma \ref{lemma6}. For any fixed $z\in (1/\sqrt{2},1)$, we interchange the integration and differentiation in the numerator in \eqref{eqn_integral_limit},   
use the factorization \eqref{E_factorization}, and obtain
\begin{equation}\label{numerator_as_sum_of_Q}
\partial_x^{m} \int_{-\infty}^y \E(z,x,u) e^{-u^2/2}  P(u) \d u
= \sum\limits_{j=0}^m (-1)^j \binom{m}{j} \widetilde {\mathcal Q}_{j}(z,x,y),
\end{equation}
where we have defined 
\begin{equation}\label{def_tilde_Q}
\widetilde {\mathcal Q}_{j}(z,x,y):=\zeta^{-1} \int_{-\infty}^y
\Big[ \partial_u^{j} e^{-(u-x)^2/\zeta^2} \Big]  \times \Big[ \partial_x^{m-j} e^{x^2/2-2xu/(1+z)} \Big]  \times P(u) \d u.
\end{equation}
When deriving the above formula, we used the fact that 
\[
\partial_x^{j} e^{-(u-x)^2/\zeta^2}=(-1)^j \partial_u^{j} e^{-(u-x)^2/\zeta^2}. 
\]
The interchange of integration and differentiation is justified since, for $z\in (1/\sqrt{2},1)$, we have $1/\zeta^2>2$; thus, for all $u\in \r$ the integrand in \eqref{def_tilde_Q} is dominated by $C\exp(- u^2)$ for some $C>0$, uniformly in $x$ on compact subsets of $\r$.  We note that the conditions  $P^{(i)}(y)=0$ for $i=0,1,\dots,n$ imply that $P(u)=q_1(u) (u-y)^{n+1}$ for some polynomial $q_1$, which in turn implies
\[
\Big[ \partial_x^{m-j} e^{x^2/2-2xu/(1+z)} \Big] \times  P(u)=
e^{x^2/2-2xu/(1+z)} q_2(x,u/(1+z),u) (u-y)^{n+1},
\]
where $q_2(x_1,x_2,x_3)$ is some polynomial in three variables.   The important conclusion is that $q_3(z,u)=q_2(x,u/(1+z),u)$ is of the form \eqref{form_of_q} (recall that everywhere in this proof $x$ is fixed, so dependence of the coefficients of the polynomial $q_3(z,u)$ on $x$ is not a concern). Thus, $\widetilde {\mathcal Q}_{j}(z,x,y)$ is an integral of the form \eqref{def_Q_j}, and using the results about ${\mathcal Q}_j(z,x,y)$ that we established above, we conclude that the limit  
\[
\lim\limits_{z\to 1-} \frac{\widetilde {\mathcal Q}_{j}(z,x,y)}{\E(z,x,y)} 
\]
exists and is equal to zero for all $0\leqslant j \leqslant n+1$, while for $j=n+2$ the limit exists and is finite. The desired result now follows from \eqref{numerator_as_sum_of_Q}.
\end{proof}

\begin{corollary}\label{Corollary2}
Let  $x$ and $y$ be fixed real numbers such that $x>y$. Then 
\[
\lim_{z\to 1-} \frac{\E_{\k}(z,x,y)}{\E(z,x,y)}
\]
exists and is finite. 
\end{corollary}
\begin{proof}
Using the same reasoning as in the derivation of \eqref{psi_k_n_sum_derivatives}, we conclude that 
\begin{equation}\label{M_kx_h}
{\mathcal M}_{\k,x} h(x)=e^{-l x^2/2} \sum\limits_{j=0}^l q_{j}(x) \partial_x^j h(x), 
\end{equation}
for some real polynomials $q_j(x)$, which may also depend on $\k$. 
The desired result now follows from Corollary \ref{Corollary1} and Lemma \ref{lemma6}. 
\end{proof}

\noindent
{\bf Proof of Theorem \ref{thm_main}:}
We denote by $\r[x]$ the set of polynomials in the $x$-variable with real coefficients, and by $\r_{n}[z,x,y]$ the set of real polynomials in the three variables $z,x,y$ whose degree in $z$ is not greater than $n$. 
An expression such as 
\[
g(z,x,y) \in H(z,x,y) \sum\limits_{i=1}^n \r[x] h_{i}(z,x,y)
\]
should be interpreted as stating that there exist polynomials $P_i \in \r[x]$ such that 
\[
g(z,x,y) = H(z,x,y) \sum\limits_{i=1}^n P_i(x) h_{i}(z,x,y)
\]
for all $z,x,y$ in the domain of $g$, $H$ and $h_i$. For example, with this notation we can express \eqref{M_kx_h} in the form 
\begin{equation}\label{item_i}
{\mathcal M}_{\k,x} h(x) \in e^{-l x^2/2}\sum\limits_{j=0}^l \r[x] \, \partial_x^j h(x). 
\end{equation}


As before, we denote $\zeta:=\sqrt{1-z^2}$ and $w:=(y-xz)/\zeta$. We collect some preliminary facts. 
\begin{itemize}
\item[(i)] For all $n\in \z_{\geqslant 0}$ 
\begin{equation}\label{item_ii}
 H_n(w)=n! \sum\limits_{j=0}^{\lfloor n/2 \rfloor} \frac{(-1)^j (2w)^{n-2j}}{j! (n-2j)!}   \in \zeta^{-n} \r_{n}[z,x,y]. 
\end{equation}
\item[(ii)] For all $n\in {\mathbb N}$ 
\begin{equation}\label{item_iii}
\partial_x^n \, {\textnormal{erf}}(w)=-(z/\zeta)^n \frac{2}{\sqrt{\pi}}  e^{-w^2} H_{n-1}(w) \in e^{-w^2} \zeta^{1-2n} \r_{2n-1}[z,x,y].
\end{equation}
\item[(iii)] For all $n,m \in \z_{\geqslant 0}$ 
\begin{equation}\label{item_iv}
\partial_x^m \Big[ e^{-w^2} H_n(w)\Big]=(z/\zeta)^m e^{-w^2} H_{n+m}(w).
\end{equation}
\end{itemize}
The expression for $H_n(w)$ in item (i) can be found  in \cite[formula 18.5.13]{NIST}. The results in items (ii) and (iii) follow from \eqref{two_integrals} and the fact that $\partial_x = (-z/\zeta) \partial_w$. 

We denote
\begin{align*}
I^{(1)}_n(z,x,y)&:=
 z^n \psi_n(x)  \frac{\sqrt{\pi}}{2} \big( 1+ 
{\textnormal{erf}}(w) \big),
\\ 
I_n^{(2)}(z,x,y)&:=-z^n e^{-w^2} \sum\limits_{i=1}^n \binom{n}{i}  (z/\zeta)^{-i}
H_{i-1}(w) \psi_{n-i}(x).
\end{align*}
By \eqref{eqn_PM}, we have 
\[
\E_{\k}(z,x,y) \in e^{-(l-1) y^2/2}  \sum\limits_{j=1}^l  \r[y] \Big[ {\mathcal M}_{\k,x} I^{(1)}_{k_j}(z,x,y)+{\mathcal M}_{\k,x} I^{(2)}_{k_j}(z,x,y)\Big].
\]
Using \eqref{item_i} we obtain
\[
{\mathcal M}_{\k,x} I^{(1)}_{k_j}(z,x,y) \in 
 z^{k_j} \frac{\sqrt{\pi}}{2} \big( 1+ 
{\textnormal{erf}}(w) \big) {\mathcal M}_{\k,x} \psi_{k_j}(x)+
 z^{k_j} e^{-lx^2/2} \sum\limits_{\substack{n\geqslant 0, m\geqslant 1 \\ m+n\leqslant l}}   \r[x] \partial_x^n \psi_{k_j}(x)\partial_x^m {\textnormal{erf}}(w).
\]
By definition of ${\mathcal M}_{\k,x}$ (see \eqref{def_M_kx}), we have ${\mathcal M}_{\k,x} \psi_{k_j}(x)\equiv 0$. We use \eqref{item_ii}, \eqref{item_iii}, and the fact that $\partial_x^n \psi_{k_j}(x)\in \exp(-x^2/2) \r[x]$ to conclude that 
\[
{\mathcal M}_{\k,x} I^{(1)}_{k_j}(z,x,y) \in 
 z^{k_j}  e^{-(l+1)x^2/2-w^2} \sum\limits_{m=1}^l   \r[x] (z/\zeta)^m H_{m-1}(w)\subset  e^{-(l+1)x^2/2-w^2} \zeta^{1-2l} \r_{2l-1+k_j}[z,x,y].
\]
In the last step, we used the fact that for all $m=1,2,\dots,l$
\[
z^{k_j+m} \zeta^{-m} H_{m-1}(w) \in 
z^{k_j+m} \zeta^{1-2m} \r_{m-1}[z,x,y] \subset \zeta^{1-2l} \r_{2l-1+k_j}[z,x,y], 
\]
which is true since $\zeta^{2l-2m}=(1-z^2)^{l-m}$. 
Similarly, using \eqref{item_ii} and \eqref{item_iv}, we conclude that 
\[
{\mathcal M}_{\k,x} I^{(2)}_{k_j}(z,x,y) 
\in z^{k_j} e^{-(l+1)x^2/2-w^2} \sum\limits_{i=1}^{k_j} \sum\limits_{m=0}^l \r[x] (z/\zeta)^{m-i} H_{m+i-1}(w) 
\subset e^{-(l+1)x^2/2-w^2} \zeta^{1-2l} \r_{2l-1+k_j}[z,x,y]. 
\]
Thus, 
\[
\E_{\k}(z,x,y) \in e^{-(l+1)x^2/2-(l-1)y^2/2-w^2} \zeta^{1-2l} \r_{2l-1+k_l}[z,x,y].
\]
Using the expression \eqref{E_factorization} and 
the above two equations we arrive at 
\[
\frac{\E_{\k}(z,x,y)}{\E(z,x,y)} \in e^{-l(x^2+y^2)/2} (1-z^2)^{1-l} \, \r_{2l-1+k_l}[z,x,y]. 
\]
Thus, we have established that
\begin{equation}\label{R_and_P}
R(z,x,y):=e^{l(x^2+y^2)/2} \frac{\E_{\k}(z,x,y)}{\E(z,x,y)}=\frac{P(z,x,y)}{(1-z^2)^{l-1}},
\end{equation}
where $P(z,x,y)$ is a real polynomial in three variables, whose degree in the $z$-variable is at most $2l-1+k_l$.  We see that $R(z,x,y)$ is a rational function of $z$ with possible poles only at $z=\pm 1$. Now assume that $x, y$ are real and $x>y$.  Letting $z\to 1-$ and applying Corollary \ref{Corollary2}, we  conclude that  $R(z,x,y)$ has a finite limit at $z=1$. This implies that, when $x, y$ are real and $x>y$, the polynomial $P(z,x,y)$ is divisible by $(1-z)^{l-1}$. This divisibility property holds for all $x,y\in \c$. Indeed, consider the Taylor coefficients of the numerator  $P(z,x,y)$ at $z=1$. These Taylor coefficients are polynomials in $x,y$ and the fact that $P(z,x,y)$ is divisible by $(1-z)^{l-1}$ implies that the first $l-1$ coefficients are equal to zero. It is clear that if a polynomial in two variables $x,y$ is zero for all real values $x>y$, then it must be equal to zero for all $x,y\in \c$. 

By the symmetry conditions \eqref{E_k_symmetries}, the rational function $R$ satisfies $R(z,x,y)=(-1)^{\kappa} R(-z,-x,y)$, which implies that it also has a finite limit as $z\to -1$. Thus, the polynomial $P(z,x,y)$ is divisible by $(1-z^2)^{l-1}$. Since $(1-z^2)^{l-1}$ has degree $2l-2$, and the numerator $P(z,x,y)$ has degree in $z$ not greater than $2l-1+k_l$, we conclude that the function $R(z,x,y)$ must be a polynomial  in the three variables $z,x,y$, with degree in $z$ not greater than $k_l+1$. 

Thus, we have established that there exists a polynomial $R(z,x,y)$, with degree in $z$ not greater than $k_l+1$, such that  
\[
\E_{\k}(z,x,y)=e^{-l(x^2+y^2)/2}\, \E(z,x,y)  R(z,x,y),
\]
which is equivalent to the desired result \eqref{eqn_main} by \eqref{mathcal_E_E}. 
\qed

\vspace{0.2cm}
\noindent
{\bf Remark 2}: One possible way to prove Conjecture 2 is to identify the coefficient of the highest power of $z$ in the polynomial $P(z,x,y)$ in \eqref{R_and_P} (namely, the coefficient of $z^{2l-1+k_l}$). Indeed, after dividing $P(z,x,y)$ by $(1-z^2)^{l-1}$, this would give $Q^{\k}_{k_l+1}(x,y)$ -- the coefficient of $z^{k_l+1}$ in the polynomial function $R(z,x,y)$. Thus, in order to prove Conjecture 2, one would need to follow the steps of the proof of Theorem \ref{thm_main} and keep track of the coefficient of the highest power of $z$ in every term. However, the resulting expressions become complicated very quickly, and we were unable to complete this computation. 
\vspace{0.2cm}

\noindent
{\bf Proof of Proposition \ref{proposition1}:}
We rewrite \eqref{eqn_main} in the form
\begin{equation}\label{generating_functions}
\sum\limits_{n \in \z_{\geqslant 0} \setminus \k}  z^n \frac{H_{\k,n}(x) H_{\k,n}(y)}{2^{n+l} n! \prod\limits_{j=1}^l (n-k_j)}= \sum\limits_{n\geqslant 0} z^n \frac{H_n(x) H_n(y)}{2^n n!} \sum_{m=0}^{k_l+1} z^m Q_m^{\k}(x,y). 
\end{equation}
Comparing the coefficient of $z^0$, we obtain
\[
Q_0^{\k}(x,y)= \frac{H_{\k,0}(x) H_{\k,0}(y)}{(-2)^{l}  \prod\limits_{j=1}^l k_j}. 
\]
Since $H_{\k,0}(x)=(-2)^l \Big[\prod_{j=1}^l k_j\Big] H_{\k-1}(x)$, which follows from \cite[Lemma 3.1]{Ullate_2018}, we obtain formula \eqref{eqn_Qk0} for  $Q_0^{\k}(x,y)$. 

Comparing the coefficients of $z^n$ in \eqref{generating_functions} gives the recursion identity \eqref{Q_m_recursion}. The parity conditions for $Q_n^{\k}(x,y)$ follow from \eqref{eqn_main} and Lemma \ref{lemma1}(iii). 
\qed


\noindent
{\bf Proof of Proposition \ref{proposition2}:}
The proof is based on the following short-time expansion of the heat kernel. This result can be found in  \cite[Section 5.9]{Avramidi2015}, see also  \cite{Iliev2005,Vassilevich2003}. Let $p(t,x,y)$ be the heat kernel of the operator  $L=-\partial_x^2+V(x)$, where $V : \r \mapsto \r$ is a smooth function. As $t\to 0^+$, we have the asymptotic expansion 
\begin{equation}\label{Hadamard_expansion}
p(t,x,y)= \frac{1}{\sqrt{4\pi t}} 
\exp\Big( -\frac{(y-x)^2}{4t} \Big) 
\big[ 1 + u_1(x,y) t + u_2(x,y) t^2+O(t^3) \big],
\end{equation}
where
\begin{equation}\label{eqn_u1}
u_1(x,y)=-\int_0^1 V(x+s(y-x)) \d s, 
\end{equation}
and 
\begin{equation}\label{eqn_u2}
u_2(x,y)=\frac{1}{2} u_1(x,y)^2 - \int_0^1 
s(1-s) V''(x+s(y-x)) \d s.
\end{equation}
The heat kernel for the operator ${\mathcal L}^{(1)}=-\partial_x^2+x^2$ is 
\[
p^{(1)}(t,x,y)=\frac{ e^{-t}}{\sqrt{\pi}} \E(e^{-2t},x,y),
\]
since ${\mathcal L}^{(1)} \psi_n=(2n+1) \psi_n$ and the functions $\psi_n$ form an orthogonal basis for $L_2(\r,\d x)$ with $\lVert \psi_n \rVert^2=\sqrt{\pi} 2^n n!$. 
Let $\omega(x):=\partial_x^2 \ln(H_{\k}(x))$. From \eqref{L_heat_kernel} we find that the heat kernel for the operator 
\[
{\mathcal L}^{(2)}=-\partial_x^2+U_{\k}(x)+2l=-\partial_x^2+x^2-2\omega(x)+2l, 
\] 
 is 
\[
p^{(2)}(t,x,y)=\frac{e^{-t} \E_{\k}(e^{-2t},x,y)}{\sqrt{\pi} \psi_{\k}(x) \psi_{\k}(y)}.
\] 
From \eqref{eqn_main} and \eqref{mathcal_E_E} we conclude that 
\begin{equation}\label{p_ratio1}
\frac{p^{(2)}(t,x,y)}{p^{(1)}(t,x,y)}=\frac{1}{H_{\k}(x) H_{\k}(y)}
\sum\limits_{n=0}^{k_l+1} e^{-2nt} Q_n^{\k}(x,y). 
\end{equation}
Let $\alpha_1(x,y)$ and $\alpha_2(x,y)$ (respectively, $\beta_1(x,y)$ and $\beta_2(x,y)$) be the coefficients $u_1(x,y)$ and $u_2(x,y)$ defined by equations \eqref{eqn_u1} and \eqref{eqn_u2} with $V(x)=x^2+2l-2\omega(x)$ (respectively, $V(x)=x^2$). For simplicity, we write $\alpha_i=\alpha_i(x,y)$ and similarly for $\beta_i$. According to \eqref{Hadamard_expansion}, we have 
\begin{equation}\label{p_ratio2}
\frac{p^{(2)}(t,x,y)}{p^{(1)}(t,x,y)}=
\frac{1+\alpha_1 t+\alpha_2 t^2+O(t^3)}{1+\beta_1 t + \beta_2 t^2 + O(t^3)}.
\end{equation}

We compute $\beta_i$ using formulas \eqref{eqn_u1} and \eqref{eqn_u2} with $V(x)=x^2$: 
\[
\beta_1=-\frac{1}{3} (x^2+xy+y^2), \;\;\; \beta_2=\frac{1}{2} \beta_1^2 - \frac{1}{3}. 
\]

Next, we use the fact that $\omega(x)=\partial_x^2 \ln(H_{\k}(x))$ and compute 
\begin{equation*}
\Omega(x,y):=\int_0^1 \omega(x+s(y-x)) \d s=\frac{1}{y-x} \Big( \frac{H'_{\k}(y)}{H_{\k}(y)}
- \frac{H'_{\k}(x)}{H_{\k}(x)} \Big). 
\end{equation*}
After integrating by parts twice, we arrive at
\begin{equation*}
\Theta(x,y):=\int_0^1 
s(1-s)\omega''(x+s(y-x)) \d s=
\frac{1}{(y-x)^2} \Big[\omega(x)+\omega(y)-2\Omega(x,y)\Big].
\end{equation*}
Again, in what follows, we will write $\Omega=\Omega(x,y)$ and $\Theta=\Theta(x,y)$. 
Using formulas \eqref{eqn_u1} and \eqref{eqn_u2} with $V(x)=x^2+2l-2\omega(x)$ we find  
\[
\alpha_1=\beta_1-2l+2 \Omega, \;\;\;
\alpha_2=\frac{1}{2} \alpha_1^2 - \frac{1}{3} + 2 \Theta. 
\]


We check that 
\[
\frac{1+\alpha_1 t+\alpha_2 t^2+O(t^3)}{1+\beta_1 t + \beta_2 t^2 + O(t^3)}=
1+\gamma_1 t + \gamma_2 t^2 + O(t^3),
\]
where 
\[
\gamma_1=\alpha_1-\beta_1=2 (\Omega-l),
\]
and 
\[
\gamma_2=\alpha_2-\beta_2-\gamma_1 \beta_1 = 2 (\Omega-l)^2 + 2\Theta.
\]
Combining the above formulas with \eqref{p_ratio1} and \eqref{p_ratio2}, we arrive at the following result: for all $x,y\in\r$, as $t\to 0^+$,
\[
\frac{1}{H_{\k}(x) H_{\k}(y)}
\sum\limits_{n=0}^{k_l+1} e^{-2nt} Q_n^{\k}(x,y)=1+2 (\Omega(x,y)-l) t + 2 \Big[ (\Omega(x,y)-l)^2 + \Theta(x,y)\Big] t^2 + O(t^3).
\]
Comparing the first three coefficients in the Taylor expansions at $t=0$ on both sides of the above equation gives us the desired formulas  \eqref{sum1}, \eqref{sum2} and \eqref{sum3}. 
\qed

\vspace{0.2cm}
\noindent
{\bf Remark 3:} Note that the asymptotic expansion \eqref{p_ratio2} implies the assertion of Corollary \ref{Corollary2}, which was used crucially in the proof of Theorem \ref{thm_main}. Thus, when $\k$ is a Krein-Adler sequence, Corollary \ref{Corollary2} is a simple consequence of the short-time heat-kernel expansion.
\vspace{0.2cm}


\noindent
{\bf Proof of Proposition \ref{proposition3}:}
Let $Q_n^k(x,y)$ be the polynomials defined in \eqref{Q_n_l=1}. Since in this case $\k=(k)$, we will write $Q_n^k(x,y)$ instead of $Q_n^{\k}(x,y)$. Similarly, we will write $H_{k,n}(x)=\wr[H_k(x),H_n(x)]$ instead of $H_{\k,n}(x)$.  
We will show that $L_n^k(x,y)=R_n^k(x,y)$ for all $k\geqslant 0$, $0\leqslant n \leqslant k+1$ and all $x,y\in \c$, where 
\[
L_n^k(x,y):=\sum_{m=0}^{n}
\frac{H_m(x)H_m(y)}{2^m m!}
Q_{n-m}^k(x,y), \;\;\;
R_n^k(x,y):=
\mathbf 1_{\{n\ne k\}}
\frac{H_{k,n}(x)H_{k,n}(y)}
{2^{n+1}n!(n-k)}.
\]
According to Proposition \ref{proposition1}, this implies that $Q_n^k(x,y)$ are the polynomials appearing in \eqref{eqn_main}. 

We denote 
\[
h_m(x,y):=H_m(x)H_m(y).
\] 
To simplify notation, we  write $h_m=h_m(x,y)$ and similarly for $Q_n^k$, $L_n^k$, and $R_n^k$. From equation \eqref{Q_n_l=1}, it follows that for $n\geqslant 1$ and $k\geqslant 1$,
\[
Q_n^k={\mathbf 1}_{\{n\geqslant 1\}} 2k Q_{n-1}^{k-1} + 
{\mathbf 1}_{\{n=1\}} h_k-2k {\mathbf 1}_{\{n=0\}} h_{k-1}. 
\]
Thus, for $n\geqslant 1$ and $k\geqslant 1$,
\[
L_n^k=\sum_{m=0}^{n}
\bigg[ \frac{h_m}{2^m m!}\bigg]
Q_{n-m}^k=2 k \sum_{m=0}^{n-1}
\bigg[ 
\frac{h_m}{2^m m!}\bigg]
Q_{n-1-m}^{k-1}+\bigg[\frac{h_{n-1}}{2^{n-1} (n-1)!}\bigg] h_k
-2k \bigg[\frac{h_{n}}{2^{n} n!}\bigg] h_{k-1},
\]
and we obtain the recurrence relation 
\begin{equation}\label{L_recursion}
L_n^k=2 k L_{n-1}^{k-1}+\frac{nh_k h_{n-1}-k h_{k-1} h_n}{2^{n-1} n!}.
\end{equation}

Next, using the facts 
\[
H_n'(x)=2n H_{n-1}(x)=2x H_n(x)-H_{n+1}(x),
\]
we check that 
\begin{align*}
&H_{k,n}(x)=2n H_k(x) H_{n-1}(x)-2k H_{k-1}(x)H_n(x), \\
&H_{k-1,n-1}(x)=H_{k}(x)H_{n-1}(x) -H_{k-1}(x) H_{n}(x). 
\end{align*}
Using the above two formulas and a straightforward but  tedious calculation, we verify that 
\[
H_{k,n}(x)H_{k,n}(y)=4n(n-k)h_k h_{n-1} - 4 k (n-k) h_{k-1} h_n+4 n k H_{k-1,n-1}(x) H_{k-1,n-1}(y).
\]
When $n\geqslant 1$ and $n\neq k$, we divide both sides of the above identity 
by $2^{n+1} n! (n-k)$ and obtain  
\begin{equation}\label{R_recursion}
R_n^k=2 k R_{n-1}^{k-1}+\frac{nh_k h_{n-1}-k h_{k-1} h_n}{2^{n-1} n!}. 
\end{equation}
We check that the above identity also holds for $n=k$, since in this case $R_n^k=R_{n-1}^{k-1}=nh_k h_{n-1}-k h_{k-1} h_n=0$.

We see from formulas \eqref{L_recursion} and \eqref{R_recursion} that $L_n^k$ and $R_n^k$ satisfy the same recurrence relation, which lowers both indices $n$ and $k$ by one. We want to prove the identity $L_n^k=R_n^k$ for $k\geqslant 0$ and $0\leqslant n \leqslant k+1$. After applying the recurrence formulas \eqref{L_recursion} and \eqref{R_recursion}  $M=\min(n,k)$ times, we arrive at either the case $n=0$, $k\geqslant 0$ or $n=1$, $k=0$. In other words, if the identity $L_n^k=R_n^k$ holds for $n=0$, $k\geqslant 0$ and  $n=1$, $k=0$, then it holds for all $k\geqslant 0$ and $n=0,1,\dots,k+1$. The fact that the identity $L_n^k=R_n^k$ holds in these boundary cases is easily verified directly using the formulas 
\[
H_{k,0}(x)=-2k H_{k-1}(x), \;\;\; H_{0,n}(x)=2n H_{n-1}(x),
\]
and the definition \eqref{Q_n_l=1} of $Q_n^k(x,y)$. 
\qed

% 
%****************************************************************************************************************
%****************************************************************************************************************
%****************************************************************************************************************e


%\bibliographystyle{abbrv}
%\bibliography{references}

\begin{thebibliography}{10}

\bibitem{Adler1994}
V.~E. Adler.
\newblock A modification of {Crum}'s method.
\newblock {\em Theoretical and Mathematical Physics}, 101(3):1381--1386, 1994.
\newblock \url{https://doi.org/10.1007/BF01035458}.

\bibitem{Avramidi2015}
I.~G. Avramidi.
\newblock {\em {Heat Kernel Method and its Applications}}.
\newblock Birkh{\"a}user, Cham, 2015.
\newblock \url{https://doi.org/10.1007/978-3-319-26266-6}.

\bibitem{Bakry2014}
D.~Bakry, I.~Gentil, and M.~Ledoux.
\newblock {\em {Analysis and Geometry of Markov Diffusion Operators}}.
\newblock Springer Cham, 2014.
\newblock \url{https://doi.org/10.1007/978-3-319-00227-9}.

\bibitem{Bonneux2018}
N.~Bonneux and M.~Stevens.
\newblock Recurrence relations for {Wronskian} {Hermite} polynomials.
\newblock {\em SIGMA. Symmetry, Integrability and Geometry: Methods and
  Applications}, 14:048, 2018.
\newblock \url{https://doi.org/10.3842/SIGMA.2018.048}.

\bibitem{NIST}
{\it NIST Digital Library of Mathematical Functions}.
\newblock \url{https://dlmf.nist.gov/}, Version 1.2.6, 2026-03-15.
\newblock F.~W.~J. Olver, A.~B. {Olde Daalhuis}, D.~W. Lozier, B.~I. Schneider,
  R.~F. Boisvert, C.~W. Clark, B.~R. Miller, B.~V. Saunders, H.~S. Cohl, and
  M.~A. McClain, eds.

\bibitem{Felder2012}
G.~Felder, A.~D. Hemery, and A.~P. Veselov.
\newblock Zeros of {Wronskians} of {Hermite} polynomials and {Young} diagrams.
\newblock {\em Physica D: Nonlinear Phenomena}, 241(23):2131--2137, 2012.
\newblock \url{https://doi.org/10.1016/j.physd.2012.08.008}.

\bibitem{Foata1978}
D.~Foata.
\newblock A combinatorial proof of the {Mehler} formula.
\newblock {\em Journal of Combinatorial Theory, Series A}, 24(3):367--376,
  1978.
\newblock \url{https://doi.org/10.1016/0097-3165(78)90066-3}.

\bibitem{Gomez-Ullate_2014}
D.~G\'omez-Ullate, Y.~Grandati, and R.~Milson.
\newblock Rational extensions of the quantum harmonic oscillator and
  exceptional {H}ermite polynomials.
\newblock {\em Journal of Physics A: Mathematical and Theoretical},
  47(1):015203, 2014.
\newblock \url{https://doi.org/10.1088/1751-8113/47/1/015203}.

\bibitem{Ullate_2018}
D.~G\'omez-Ullate, Y.~Grandati, and R.~Milson.
\newblock {Durfee rectangles and pseudo-Wronskian equivalences for Hermite
  polynomials}.
\newblock {\em Studies in Applied Mathematics}, 141(4):596--625, 2018.
\newblock \url{https://doi.org/10.1111/sapm.12225}.

\bibitem{Gomez2021}
D.~G{\'o}mez-Ullate, Y.~Grandati, and R.~Milson.
\newblock Complete classification of rational solutions of
  {$A_{2n}$}-{Painlev{\'e}} systems.
\newblock {\em Advances in Mathematics}, 385:107770, 2021.
\newblock \url{https://doi.org/10.1016/j.aim.2021.107770}.

\bibitem{Gomez2016}
D.~G{\'o}mez-Ullate, A.~Kasman, A.~B.~J. Kuijlaars, and R.~Milson.
\newblock Recurrence relations for exceptional {Hermite} polynomials.
\newblock {\em Journal of Approximation Theory}, 204:1--16, 2016.
\newblock \url{https://doi.org/10.1016/j.jat.2015.12.003}.

\bibitem{Jeffrey2007}
I.~S. Gradshteyn and I.~M. Ryzhik.
\newblock {\em Table of Integrals, Series, and Products}.
\newblock Academic Press, 7 edition, 2007.
\newblock edited by A. Jeffrey and D. Zwillinger.

\bibitem{Iliev2005}
P.~Iliev.
\newblock On the heat kernel and the {Korteweg--de Vries} hierarchy.
\newblock {\em Annales de l'Institut Fourier}, 55(6):2117--2127, 2005.
\newblock \url{https://doi.org/10.5802/aif.2154}.

\bibitem{Namias1980}
V.~Namias.
\newblock The fractional order {Fourier} transform and its application to
  quantum mechanics.
\newblock {\em IMA Journal of Applied Mathematics}, 25(3):241--265, 1980.
\newblock \url{https://doi.org/10.1093/imamat/25.3.241}.

\bibitem{Oblomkov1999}
A.~A. Oblomkov.
\newblock Monodromy-free {Schr{\"o}dinger} operators with quadratically
  increasing potentials.
\newblock {\em Theoretical and Mathematical Physics}, 121(3):1574--1584, 1999.
\newblock \url{https://doi.org/10.1007/BF02557204}.

\bibitem{Pupasov_2015}
A.~M. Pupasov-Maksimov.
\newblock {Propagators of isochronous an-harmonic oscillators and Mehler
  formula for the exceptional Hermite polynomials}.
\newblock {\em Annals of Physics}, 363:122--135, 2015.
\newblock \url{https://doi.org/10.1016/j.aop.2015.09.021}.

\bibitem{Pupasov_2007}
A.~M. Pupasov-Maksimov, B.~F. Samsonov, and U.~G\"unther.
\newblock {Exact propagators for SUSY partners}.
\newblock {\em Journal of Physics A: Mathematical and Theoretical},
  40(34):10557, 2007.
\newblock \url{https://dx.doi.org/10.1088/1751-8113/40/34/013}.

\bibitem{Szego}
G.~Szeg\"o.
\newblock {\em {Orthogonal polynomials}}.
\newblock Colloquium Publications, Volume XXIII. American Mathematical Society,
  Providence, Rhode Island, 1939.

\bibitem{Vassilevich2003}
D.~V. Vassilevich.
\newblock Heat kernel expansion: user's manual.
\newblock {\em Physics Reports}, 388(5-6):279--360, 2003.
\newblock \url{https://doi.org/10.1016/j.physrep.2003.09.002}.

\end{thebibliography}


\end{document}


